\section{The Software-as-a-Graph (SaG) Architectural Model}
\label{sec:3}

This section formalizes the Software-as-a-Graph multigraph representation (Section~\ref{sec:3.1}), the QoS-aware weighting and logical dependency derivation rules (Section~\ref{sec:3.2}), the dual graph views (Section~\ref{sec:3.3}), and the typed node feature encodings (Section~\ref{sec:3.4}).

\subsection{Formal Multigraph Definition}
\label{sec:3.1}

A complex distributed software system is formally modeled as a typed, weighted, directed multigraph:
\begin{equation}
\mathcal{G} = (V, E, \tau_V, \tau_E, w_V, w_E)
\end{equation}
where:
\begin{itemize}
    \item $V$ is the set of system entities, partitioned into five disjoint entity types:
    \begin{equation}
    V = V_{\text{app}} \cup V_{\text{broker}} \cup V_{\text{topic}} \cup V_{\text{node}} \cup V_{\text{lib}}
    \end{equation}
    \item $E$ is the set of directed edges connecting entities.
    \item $\tau_V: V \to \mathcal{T}_V$ and $\tau_E: E \to \mathcal{T}_E$ are typing functions assigning node and edge categories.
    \item $w_V: V \to [0, 1]$ and $w_E: E \to [0, 1]$ are weighting functions representing entity criticality and connection strength.
\end{itemize}

Table~\ref{tab:1} summarizes the five entity types and six structural edge types formalized in the SaG model, along with their semantics and representative concrete distributed-system implementations.

\begin{table}[htbp]
\centering
\small
\caption{Entity and structural edge types in the SaG model.}
\label{tab:1}
\resizebox{\linewidth}{!}{%
\begin{tabular}{lll}
\toprule
\textbf{Entity Type ($\mathcal{T}_V$)} & \textbf{Architectural Role} & \textbf{Concrete System Examples} \\
\midrule
\textbf{Application} ($V_{\text{app}}$) & Autonomous process producing/consuming messages & ROS~2 node, Kafka microservice, MQTT client \\
\textbf{Broker} ($V_{\text{broker}}$) & Message routing and queuing intermediary & RabbitMQ exchange, Mosquitto, EMQX broker \\
\textbf{Topic} ($V_{\text{topic}}$) & Named logical communication channel & \texttt{/sensor/lidar}, \texttt{orders.payment.completed} \\
\textbf{Node} ($V_{\text{node}}$) & Physical host or virtualized execution environment & Bare-metal server, Kubernetes worker, Cloud VM \\
\textbf{Library} ($V_{\text{lib}}$) & Shared software package or runtime dependency & \texttt{librdkafka}, OpenCV, Protobuf runtime \\
\midrule
\textbf{Structural Edge ($\mathcal{T}_E$)} & \textbf{Direction} & \textbf{Semantic Meaning} \\
\midrule
\texttt{PUBLISHES\_TO} & App/Library $\to$ Topic & Component publishes messages to topic \\
\texttt{SUBSCRIBES\_TO} & App/Library $\to$ Topic & Component consumes messages from topic \\
\texttt{ROUTES} & Broker $\to$ Topic & Broker manages and routes topic traffic \\
\texttt{RUNS\_ON} & App/Broker $\to$ Node & Process is hosted on physical/virtual host \\
\texttt{CONNECTS\_TO} & Node $\to$ Node & Physical network link between hosts \\
\texttt{USES} & App $\to$ Library & Application links to shared library dependency \\
\bottomrule
\end{tabular}%
}
\end{table}

Application and Library entities additionally incorporate static code metrics computed via Static Code Analysis (SCA) tools (\texttt{cm\_*} attributes: lines of code, cyclomatic complexity, coupling between objects, LCOM), linking code-level fragility directly to topological analysis.

\subsection{QoS-Aware Weights and Logical Dependency Derivation}
\label{sec:3.2}

In distributed middleware, communication links vary in coupling strength based on their Quality-of-Service (QoS) contracts. For instance, a \texttt{RELIABLE} topic with \texttt{TRANSIENT\_LOCAL} durability binds communicating services substantially more tightly than a \texttt{BEST\_EFFORT} telemetry stream.

Each topic $t$ carries an intrinsic criticality weight $w(t) \in [0, 1]$ combining its declared QoS semantics with two runtime-stress modulators: payload size and publication frequency:
\begin{equation}
\label{eq:3}
w(t) = \beta \cdot \text{QoS}(t) + \alpha \cdot \text{SizeNorm}(t) + \psi \cdot \text{FreqNorm}(t),
\quad (\beta, \alpha, \psi) = (0.75,\, 0.15,\, 0.10)
\end{equation}
where the QoS term is an AHP-weighted aggregate of the declared contract:
\begin{equation}
\text{QoS}(t) = w_{\text{rel}} \cdot q_{\text{rel}} + w_{\text{dur}} \cdot q_{\text{dur}} + w_{\text{prio}} \cdot q_{\text{prio}},
\quad (w_{\text{rel}}, w_{\text{dur}}, w_{\text{prio}}) = (0.24,\, 0.62,\, 0.14)
\end{equation}
Here, $q_{\text{rel}}, q_{\text{dur}}, q_{\text{prio}} \in [0, 1]$ represent normalized reliability, durability, and transport-priority scores. Durability dominates because it governs whether data persists across restarts and network partitions. Reliability and transport priority both govern in-flight delivery quality, with reliability receiving higher weight because unconditional delivery guarantees precede message scheduling. The sub-weight vector is the geometric-mean priority vector of an independently stated Saaty pairwise-comparison matrix, printed with its consistency computation in Supplementary Section~S4; $CR = 0.016$, small but non-zero. The modulators are logarithmically compressed and clamped to $[0, 1]$: $\text{SizeNorm}(t) = \min(1.0, \log_2(1 + \text{bytes})/20)$ (a 1~MiB design envelope, representing the practical DDS sample ceiling before RTPS fragmentation dominates) and $\text{FreqNorm}(t) = \min(1.0, \log_{10}(1 + \text{Hz})/3)$. The final weight $w(t)$ is clamped to $[0.01, 1]$, ensuring that best-effort edges remain visible to graph traversals. Every structural communication edge incident on $t$ (\texttt{PUBLISHES\_TO}, \texttt{SUBSCRIBES\_TO}, \texttt{ROUTES}) inherits $w_E(e) = w(t)$ along with the topic's QoS vector.

The outer split $(\beta, \alpha, \psi)$ is a declared convex combination. Sweeping it over the full simplex changes the induced ordering of $w(t)$ by at most $\rho = 0.081$ and downstream rank correlation by at most $0.031$, so it is a documented convention rather than a tuned parameter (Supplementary Section~S1).

\subsubsection{Logical Dependency Projection (\texttt{DEPENDS\_ON})}

Structural edges capture explicit deployment connections but omit implicit runtime dependencies. For example, a subscriber depends upon a publisher, yet no direct edge connects them in pub-sub architectures. We therefore derive a single unified semantic relation, \texttt{DEPENDS\_ON}, directed from \emph{dependent} to \emph{dependency} (``if target fails, source is impacted''), according to the six projection rules detailed in Table~\ref{tab:2}:

\begin{table}[htbp]
\centering
\small
\caption{The six \texttt{DEPENDS\_ON} logical dependency projection rules.}
\label{tab:2}
\resizebox{\linewidth}{!}{%
\begin{tabular}{clll}
\toprule
\textbf{Rule} & \textbf{Dependency Category} & \textbf{Structural Pattern ($\text{Dependent} \to \text{Dependency}$)} & \textbf{Derived Weight ($w$)} \\
\midrule
\textbf{1} & \texttt{app\_to\_app} & Subscriber $\to$ Publisher (via shared Topic, incl.\ transitive \texttt{USES}) & $1 - \prod_{t \in T}(1 - w(t))$ \\
\textbf{2} & \texttt{app\_to\_broker} & Publisher/Subscriber $\to$ Broker routing its topics & $1 - \prod_{t \in T}(1 - w(t))$ \\
\textbf{3} & \texttt{node\_to\_node} & Host $\to$ Host (lifted from inter-host app dependencies) & Lifted $\max w$ \\
\textbf{4} & \texttt{node\_to\_broker} & Host $\to$ Broker (lifted from hosted app dependencies) & Lifted $\max w$ \\
\textbf{5} & \texttt{app\_to\_lib} & Application $\to$ Shared Library it \texttt{USES} & $H(w_V(\text{app}), w_V(\text{lib}))$ \\
\textbf{6} & \texttt{broker\_to\_broker} & Broker $\leftrightarrow$ Broker (shared physical fault-domain colocation, symmetric) & $w_V(\text{node})$ \\
\bottomrule
\end{tabular}%
}
\end{table}

Rules 1 and 2 aggregate the set of topics $T$ connecting a component pair using a probabilistic union rather than a maximum \cite{pearl1988probabilistic,beliakov2007aggregation,yager1988owa}. This guarantees that additional parallel failure vectors increase coupling monotonically while keeping $w \in (0, 1]$. Rule 5 applies the harmonic mean $H(x, y) = 2xy/(x+y)$ \cite{hardy1952inequalities} to combine the consuming Application's and the shared Library's vertex weights, balancing caller and dependency criticality. Rules 3 and 4 assign the maximum weight among component-level dependencies crossing the host boundary.

\subsubsection{Sequential Cascades vs.\ Simultaneous Blasts}

A foundational principle of the SaG model is distinguishing two degradation modes: (1)~\textbf{Sequential Cascades (Rule 1)}, where a failed publisher starves downstream subscribers hop by hop through message queues and topic buffers; and (2)~\textbf{Simultaneous Blasts (Rule 5)}, where a crashed library or execution node causes all consuming applications and colocated brokers to fail instantaneously in a single shared-fate event. Preserving entity types and relation-specific projection rules enables SaG to model both mechanisms, whereas untyped homogeneous graphs collapse them into indistinguishable edges.

Rule 6 is the only symmetric projection rule: two brokers colocated on the same host share that host's physical failure domain (following the simultaneous-blast principle of Rule 5, with $w = w_V(\text{node})$). In production deployments, colocated brokers compete for CPU, memory, and I/O; a host crash halts all colocated instances simultaneously. Rule 6 does not model logical intra-cluster broker coupling (e.g., partition replication, quorum election, or shovel links, which do not require physical colocation). It applies in four of the eight detection benchmark scenarios (Supplementary Section~S12) and contributes only 12 directed edges. Because simulation oracles operate strictly on $G_{\text{structural}}$ (Section~\ref{sec:4.4}), Rule 6 has zero influence on ground-truth failure labels.

\subsection{Dual Graph Views and Architectural Layers}
\label{sec:3.3}

The SaG framework maintains two distinct representations: (1)~\textbf{Structural Graph ($G_{\text{structural}}$)}, the raw deployment graph containing physical relations (\texttt{PUBLISHES\_TO}, \texttt{ROUTES}, \texttt{RUNS\_ON}, \texttt{USES}) consumed exclusively by discrete-event simulators for unbiased failure injections (Section~\ref{sec:4.3}); and (2)~\textbf{Analysis Graph ($G_{\text{analysis}}$)}, the projected graph containing derived \texttt{DEPENDS\_ON} edges annotated with QoS weights and ingested SCA metrics, over which all GNN embeddings and analytical metrics are computed (see running example in Supplementary Figure~S14.1).

$G_{\text{analysis}}$ is further structured into four analytical layers (Application, Middleware, Infrastructure, and Global System), enabling evaluation of criticality at subsystem levels, consistent with hierarchical frameworks such as MIL-STD-498 \cite{dod1994milstd498}.

\subsection{Typed Node Feature Encoding}
\label{sec:3.4}

Both pathways read the same typed node properties from $G_{\text{analysis}}$: the predictive pathway (Section~\ref{sec:4}) projects them per entity type before heterogeneous message passing, and the explanation layer (Section~\ref{sec:5}) aggregates them into its quality profile. All five entity types share indices 0--17, a common block of topological metrics --- in/out degree, betweenness, closeness, reverse PageRank, clustering coefficient, articulation score and bridge load --- produced by the deterministic analysis stage whose cost is characterized in Section~\ref{sec:7.5}. All topological metrics in this block are normalized to $[0, 1]$ within each graph: degrees are normalized by $|V|-1$, betweenness and closeness follow standard network formulations, and reverse PageRank is normalized to unit sum. This within-graph normalization prevents raw graph size and component counts from dominating multi-layer perceptron projections during cross-scenario inductive transfer. Type-specific blocks extend it to between 19 and 25 dimensions: source-code metrics and the Code Quality Penalty for Applications, two reverse-\texttt{USES} blast-radius drivers for Libraries, queue capacity for Brokers, publisher/subscriber counts and ordinal QoS criticality for Topics, and CPU and memory allocation for Infrastructure Nodes. Supplementary Section~S11 gives the index-by-index schema.

That the shared block is where the graph structure lives matters for interpreting Section~\ref{sec:7}: betweenness, closeness, reverse PageRank and articulation score are already summaries of the topology, computed before any model sees the graph. A learned model is therefore not the only route from structure to a criticality score, which is what makes the closed-form baselines fair comparators rather than strawmen.
